\documentclass{article}

\usepackage{amsmath}
\usepackage{amssymb}
\usepackage{bm}
\usepackage[colorlinks=true, linkcolor=blue, citecolor=blue, urlcolor=blue]{hyperref}
\usepackage[margin=2.5cm]{geometry}
\usepackage{booktabs}
\usepackage{graphicx}
\usepackage{listings}
\usepackage{xcolor}

\lstset{
    language=Matlab,
    basicstyle=\ttfamily\small,
    keywordstyle=\color{blue},
    commentstyle=\color{green!50!black},
    stringstyle=\color{red!60!black},
    numbers=left,
    numberstyle=\tiny\color{gray},
    frame=single,
    breaklines=true,
    captionpos=b
}

\newcommand{\exampleimage}[2]{%
    \IfFileExists{#1}{%
        \includegraphics[width=\linewidth]{#1}%
    }{%
        \fbox{\parbox[c][0.28\textheight][c]{0.92\linewidth}{\centering #2}}%
    }%
}

\title{Frame-Based Solver for 3D Manipulator Arm:\\
Static Condensation with Kinematic Constraints}
\author{}
\date{}

\begin{document}

\maketitle
\tableofcontents
\bigskip

%-----------------------------------------------------------------------------
\section{Introduction}
\label{sec:intro}

The structural analysis of a multi-segment continuum manipulator arm involves
two distinct spatial scales:
\begin{itemize}
    \item \textbf{Skeletal (frame) scale}: a 1D beam model capturing the global
          deformation of each tubular segment via 6-DOF Timoshenko frame
          elements (Cowper shear correction for hollow circular sections).
    \item \textbf{Cross-sectional (solid) scale}: a full 3D hexahedral finite
          element model resolving stress distributions within each segment,
          including the inclined cut geometry at the segment joints.
\end{itemize}

Solving the full 3D solid model directly for every topology-optimisation
iteration is computationally prohibitive.  The \emph{frame-based solver}
exploits the separation of scales: the light-weight frame model is solved
first to obtain joint displacements, which are then imposed as kinematic
boundary conditions on the 3D solid model.  Static condensation eliminates
the interior DOFs analytically, yielding an efficient solution of the solid
model without iterative solvers on the full system.

%-----------------------------------------------------------------------------
\section{Frame Model}
\label{sec:frame}

\subsection{Degrees of Freedom}

Each frame node $k$ carries six DOFs assembled into the displacement vector
\begin{equation}
    \bm{q}_k = \bigl[u_x,\; u_y,\; u_z,\; \theta_x,\; \theta_y,\; \theta_z\bigr]_k^{\mathsf{T}} \in \mathbb{R}^6.
\end{equation}
For a model with $n_f$ frame nodes the global frame displacement vector is
$\bm{q}_{\text{frame}} \in \mathbb{R}^{6 n_f}$.

\subsection{Stiffness and Solution}

The Timoshenko element stiffness matrix $\bm{K}_e^{\text{frame}} \in \mathbb{R}^{12\times12}$
is assembled for each beam element connecting nodes $I$ and $J$ (see the shear
correction discussion in Section~\ref{sec:locking} for the shear parameter
$\Phi$).  After applying displacement boundary conditions (fixed base node), the frame system
\begin{equation}
    \bm{K}_{\text{frame}}\,\bm{q}_{\text{frame}} = \bm{P}_{\text{frame}}
\end{equation}
is solved directly.  In the MATLAB implementation this is handled by
\lstinline{frame_analysis.solveWeighted(ones(nFE_frame, 1))}, which assembles
and factorises $\bm{K}_{\text{frame}}$ and stores the result in
\lstinline{frame_analysis.qnodal} (shape: $n_f \times 6$).

%-----------------------------------------------------------------------------
\section{3D Solid Model}
\label{sec:solid}

\subsection{Degrees of Freedom}

Each solid mesh node $i$ carries three translational DOFs
\begin{equation}
    \bm{u}_i = \bigl[u_x,\; u_y,\; u_z\bigr]_i^{\mathsf{T}} \in \mathbb{R}^3.
\end{equation}
For a mesh with $n_s$ solid nodes the global displacement vector is
$\bm{u} \in \mathbb{R}^{3 n_s}$.

\subsection{Hexahedral Elements}

Each segment is discretised with 8-node serendipity hexahedral elements
(\texttt{ShapeFunctionL8}).  The element stiffness matrix is
\begin{equation}
    \bm{K}_e^{\text{solid}} = \int_{\Omega_e} \bm{B}^{\mathsf{T}} \bm{D} \bm{B}\,\mathrm{d}V,
\end{equation}
where $\bm{B}$ is the strain--displacement matrix and $\bm{D}$ is the
linear-elastic constitutive matrix for an isotropic material with Young's
modulus $E$ and Poisson's ratio $\nu$.

%-----------------------------------------------------------------------------
\section{Inclined Cross-Section Geometry}
\label{sec:geometry}

\subsection{Cut Rotation}

The joints between adjacent segments are \emph{inclined cuts}: the contact
surface is not perpendicular to the beam axis but tilted by angle $\alpha$
about the $y$-axis.  The rotation matrix that maps the beam-axis frame to the
cut-surface frame is
\begin{equation}
    \bm{R}_{\text{cut}} = \bm{R}_y(\alpha)^{\mathsf{T}}
    = \begin{bmatrix}
         \cos\alpha & 0 &  \sin\alpha \\
         0          & 1 &  0          \\
        -\sin\alpha & 0 &  \cos\alpha
      \end{bmatrix}^{\mathsf{T}}.
\end{equation}
Because the end-face normal is not aligned with the segment axis, a node at
radial distance $\rho$ from the axis spans an axial range of
$\rho\sin\alpha$ on the inclined face, reaching $R\sin\alpha$ at the outer
wall.  For $R = 0.14\,\mathrm{m}$ and $\alpha = 22.5^\circ$ this equals
approximately $54\,\mathrm{mm}$.

\subsection{Exact Section-Node Tracking During Mesh Generation}

Coupling node sets are registered \emph{during mesh generation} rather than
recovered by a post-hoc geometric search.  This guarantees exactness
regardless of the inclination angle and avoids tolerance tuning.

Each segment generator (\texttt{generateSegment2a}, \texttt{generateSegment2b})
works in a reference parameter space where the axial coordinate is $\zeta \in [0,1]$.
Nodes at the start and end faces are identified by their extreme axial parameter
values, before any cylindrical or rotation mapping is applied:
\begin{equation}
    \mathcal{S}_{\text{start}} = \bigl\{i : |\zeta_i - \zeta_{\min}| \leq \varepsilon_h\bigr\},
    \qquad
    \mathcal{S}_{\text{end}}   = \bigl\{i : |\zeta_i - \zeta_{\max}| \leq \varepsilon_h\bigr\},
\end{equation}
where $\varepsilon_h = 1/\texttt{mesh.tolerance}$ is the snapping granularity.
Because $\zeta$ is an integer grid parameter in the reference mesh, the start
nodes have $\zeta = 0$ and the end nodes have $\zeta = \texttt{resLen}$, so
the selection is exact.

\subsection{Coordinate Mapping After Mesh Merging}

After each call to \texttt{Mesh.merge}, the global mesh accumulates nodes from
multiple segments and duplicate nodes at shared faces are collapsed.  The
section node lists (recorded in local mesh indices) are re-mapped to global
indices by coordinate lookup:
\begin{enumerate}
    \item Round all node coordinates to the mesh snapping grid:
          $\tilde{\bm{x}} = \mathrm{round}(\bm{x} \cdot \texttt{tolerance})$.
    \item For each section node coordinate, find its row in the global
          rounded-coordinate table via \texttt{ismember}.
    \item The resulting global indices form the coupling section
          $\mathcal{S}_k^{(\text{global})}$.
\end{enumerate}
This lookup is exact (no floating-point tolerance) because both the node
coordinates and the snapping grid are identical up to rounding.

\subsection{Coupling Section Registry}

The method \texttt{registerCouplingSection} appends a record
\begin{equation}
    \bigl\{\,\texttt{frameNode}=k,\;\texttt{adjacentFrameNode}=\ell,\;
    \texttt{nodes}=\mathcal{S}_k^{(\text{global})},\;\texttt{label}\bigr\}
\end{equation}
to \texttt{obj.couplingSections} for every section registered during arm
assembly.  Interior joints accumulate two records (one from each incident
segment) that share the same frame node index; after merging, their node
sets are identical so the kinematic imposition is consistent.

%-----------------------------------------------------------------------------
\section{Rigid-Body Kinematic Constraint}
\label{sec:kinematics}

\subsection{Constraint Equation}

At each frame node $k$ (centroid $\bm{c}_k$, displacement $\bm{u}_k$,
rotation $\bm{\theta}_k$), the prescribed displacement for a solid node
$\bm{p}$ on the associated cross-section is given by the linearised
rigid-body formula:
\begin{equation}
    \boxed{\bm{u}(\bm{p}) = \bm{u}_k + \bm{\theta}_k \times (\bm{p} - \bm{c}_k)}
    \label{eq:rigid_body}
\end{equation}

\subsection{Linearisation Assumption}

Equation~\eqref{eq:rigid_body} is exact for rigid bodies undergoing
small displacements and small rotations (infinitesimal rotation
approximation).  This is the standard assumption in linear structural
mechanics and is consistent with the underlying FEM formulation.  For a
slender elastic tube the cross-section rotation is small (Euler--Bernoulli
hypothesis), so the formula provides an accurate leading-order approximation.

%-----------------------------------------------------------------------------
\section{Static Condensation}
\label{sec:condensation}

\subsection{DOF Partition}

After identifying all boundary DOFs $\mathcal{B}$ (prescribed by the frame
solution via \eqref{eq:rigid_body}) and interior DOFs $\mathcal{I}$, the
global stiffness system is partitioned as
\begin{equation}
    \begin{bmatrix} \bm{K}_{ii} & \bm{K}_{ib} \\ \bm{K}_{bi} & \bm{K}_{bb} \end{bmatrix}
    \begin{bmatrix} \bm{q}_i \\ \bm{q}_b \end{bmatrix}
    =
    \begin{bmatrix} \bm{P}_i \\ \bm{P}_b \end{bmatrix}.
\end{equation}

\subsection{Condensed Interior System}

Since the solid model is driven purely by kinematic loading ($\bm{P}_i = \bm{0}$,
no body forces or surface tractions on the interior), the interior equilibrium
reduces to
\begin{equation}
    \bm{K}_{ii}\,\bm{q}_i = -\bm{K}_{ib}\,\bm{q}_b.
    \label{eq:condensed}
\end{equation}
Equation~\eqref{eq:condensed} is a square, non-singular system (provided
$\mathcal{B}$ is a sufficient constraint set) and is solved by sparse direct
factorisation.  The boundary values $\bm{q}_b$ are set to the prescribed
values derived from \eqref{eq:rigid_body}.

\subsection{Choice of Boundary Set $\mathcal{B}$: Avoiding Kinematic Locking}
\label{sec:locking}

A naive choice of $\mathcal{B}$ — prescribing rigid-body kinematics at
\emph{every} cross-section ring in the arm (i.e., at each of the $n_f$
frame joints) — over-constrains the solid model and produces
\textbf{kinematic locking}.

The mechanism is as follows.  Consider one segment between frame nodes $k$
and $k+1$.  If both the start ring (at $k$) and the end ring (at $k+1$) are
rigidly prescribed with independently computed frame displacements, the
interior solid nodes must accommodate the mismatch between two rigid-body
fields that are in general \emph{inconsistent} with the elastic equilibrium of
the tube.  The resulting artificial compatibility stresses dominate the
interior field, completely masking the physically meaningful stress distribution.

\paragraph{Correct formulation: displacement control at root and tip only.}
The correct boundary set contains only two groups:
\begin{enumerate}
    \item \textbf{Root cross-section} (frame node 1): prescribe $\bm{u}(\bm{p}) = \bm{0}$
          (rigid clamp, consistent with the frame fixed-base condition).
    \item \textbf{Tip cross-section} (last frame node $n_f$): prescribe
          $\bm{u}(\bm{p}) = \bm{u}_{n_f} + \bm{\theta}_{n_f} \times (\bm{p} - \bm{c}_{n_f})$
          from the frame tip displacement.
\end{enumerate}
All interior cross-section rings are left \emph{free}.  The arm deforms
elastically between the two end constraints, producing smooth interior
stresses.  This is a pure \emph{displacement-controlled} problem ($\bm{P} =
\bm{0}$): no external forces are applied, and the load is transmitted entirely
through the prescribed end displacements.

\paragraph{Shear deformation and Timoshenko corrections.}
For non-slender beams the Euler--Bernoulli assumption (plane sections remain
plane, no shear deformation) over-stiffens the element.  The relevant
dimensionless parameter is
\begin{equation}
    \Phi = \frac{12 E I}{\kappa G A l^2},
\end{equation}
where $\kappa$ is the cross-section shear correction factor.  Euler--Bernoulli
overestimates the bending stiffness by a factor $(1+\Phi)$ relative to the
Timoshenko beam.  For the hollow tubular arm studied here
($E = 2\,\mathrm{GPa}$, $\nu = 0.35$, $R_o = 0.14\,\mathrm{m}$,
$R_i = 0.13\,\mathrm{m}$, $l_s = 0.15\,\mathrm{m}$) one obtains
$\Phi \approx 26$ for the end segments and $\Phi \approx 7$ for the interior
segments.  Using Euler--Bernoulli would therefore underestimate the tip
displacement by a factor of up to~27, leading to proportionally smaller
prescribed BCs on the solid and a 2--3$\times$ underestimate of stress
magnitudes.

The implementation uses the \textbf{Timoshenko stiffness matrix}:
\begin{align}
    k_{2,2} &= \frac{12 E J_z}{l^3(1+\Phi_z)}, &
    k_{6,6} &= \frac{(4+\Phi_z) E J_z}{l(1+\Phi_z)}, &
    k_{6,2} &= \frac{6 E J_z}{l^2(1+\Phi_z)}, \\
    k_{12,6} &= \frac{(2-\Phi_z) E J_z}{l(1+\Phi_z)}, &
    \multispan2{\text{and symmetric counterparts in } x\text{-}z \text{ plane.}}
\end{align}
The shear correction factor for a hollow circular cross-section is given by
the Cowper (1966) formula:
\begin{equation}
    \kappa = \frac{6(1+\nu)(1+m^2)^2}{(7+6\nu)(1+m^2)^2 + (20+12\nu)m^2},
    \qquad m = R_i/R_o.
\end{equation}
For $\nu = 0.35$ and $m = 13/14 \approx 0.929$ this gives $\kappa \approx 0.53$.

%-----------------------------------------------------------------------------
\section{Algorithm Summary}
\label{sec:algorithm}

\begin{enumerate}
    \item \textbf{Frame solve}: assemble and solve
          $\bm{K}_{\text{frame}}\,\bm{q}_{\text{frame}} = \bm{P}_{\text{frame}}$
          to obtain $\bm{q}_{\text{frame}}$ (joint translations and rotations).

    \item \textbf{Coupling section retrieval}: call
          \texttt{buildCouplingSections()} to retrieve the pre-registered
          per-section node lists that were built exactly during arm mesh
          assembly (Section~\ref{sec:geometry}).  No post-hoc geometric
          search is performed at solve time.

    \item \textbf{Kinematic imposition}: apply rigid-body kinematic BCs
          \eqref{eq:rigid_body} only at the \textbf{root} (frame node 1,
          $\bm{u}=\bm{0}$) and the \textbf{tip} (frame node $n_f$,
          frame-computed displacement).  Interior joint cross-sections are
          left free to avoid kinematic locking (Section~\ref{sec:locking}).

    \item \textbf{Global stiffness assembly}: assemble
          $\bm{K}_{\text{solid}}(x) = \sum_e x_e \bm{K}_e^{\text{solid}}$
          for density field $x$ (used in topology optimisation).

    \item \textbf{Static condensation solve}: use \texttt{LinearEquationsSystem.solvePq}
          to solve \eqref{eq:condensed} with the combined support set.

    \item \textbf{Result recovery}: restore prescribed DOF values, store
          $\bm{u}$ in \texttt{qnodal\_solid}, and invoke
          \texttt{computeElementResults} for stresses and derived quantities.

    \item \textbf{Frame force recovery}: call \texttt{frameElem.computeResults}
          to extract local internal forces from the frame solution.
\end{enumerate}

%-----------------------------------------------------------------------------
\section{MATLAB Implementation Notes}
\label{sec:matlab}

\subsection{\texttt{solvePq} Interface}

\texttt{LinearEquationsSystem.solvePq(K\_vals, P, q0)} solves the
system $\bm{K}\bm{q} = \bm{P}$ subject to support constraints, where
\texttt{q0} provides the prescribed values at supported DOFs.  The method
internally partitions the system, solves the reduced system, and returns
the full displacement vector.  After the call, the prescribed DOF values
must be restored explicitly since \texttt{solvePq} returns zero at
supported DOFs:
\begin{lstlisting}[caption={Restoring prescribed boundary values after solvePq}]
q_fem_sol = solver_sc.solvePq(K_vals, P_zero, q0_fem);
q_fem_sol(solver_sc.supdofs) = q0_fem(solver_sc.supdofs);
\end{lstlisting}

\subsection{FEM Vector Ordering}

The solid FEM uses DOF ordering interleaved by node: for node $i$ with 3 DOFs
$(u_x, u_y, u_z)$, the global DOF indices are $3(i-1)+1$, $3(i-1)+2$,
$3(i-1)+3$.  The helper \texttt{fromFEMVector} reshapes the flat FEM vector
back to the $n_s \times 3$ nodal matrix stored in \texttt{qnodal}.

\subsection{DOF Conventions}

\begin{table}[h]
\centering
\caption{DOF convention summary}
\begin{tabular}{@{}lll@{}}
\toprule
Model       & DOFs per node & Ordering \\
\midrule
Frame       & 6             & $[u_x, u_y, u_z, \theta_x, \theta_y, \theta_z]$ \\
3D Solid    & 3             & $[u_x, u_y, u_z]$ \\
\bottomrule
\end{tabular}
\end{table}

\subsection{Support Matrix Layout}

The \texttt{analysis.supports} matrix has shape $n_s \times 3$ (one row per
solid node, one column per translational DOF).  When merging kinematic
constraints from the frame solution with existing boundary conditions, a
logical OR is applied:
\begin{lstlisting}[caption={Combining original and kinematic supports}]
sup_combined = logical(obj.analysis.supports) | supports_new;
sup_fem      = reshape(sup_combined', [], 1);
\end{lstlisting}
Note the transpose before reshape: MATLAB stores matrices column-major, so
transposing ensures the DOF ordering matches the stiffness matrix assembly.

%-----------------------------------------------------------------------------
\section{Illustrative Coupled Examples}
\label{sec:examples}

The script \texttt{examples/ManipulatorArm/coupledExamples.m} generates four
illustrative coupled frame-solid cases and exports their images directly to
the \texttt{docs/} directory.  Each case uses the same tubular-arm template,
solves the frame-constrained solid model with \texttt{frameBasedSolver(1)},
solves the corresponding direct 3D solid model with \texttt{compute(1)}, and
writes \emph{three} comparison figures per configuration:
\begin{itemize}
    \item a deformed-mesh view (coupled on the left, direct on the right),
    \item a Huber--Mises stress comparison with a \textbf{shared} color scale
          $[\sigma_{\min},\,\sigma_{\max}]$ computed from both models combined
          — this reveals absolute magnitude differences but may suppress
          distribution details in the lower-stress panel,
    \item a Huber--Mises stress comparison with \textbf{independent} per-panel
          color scales $[0, \mathrm{HM}_{\max}^{\text{coupled}}]$ and
          $[0, \mathrm{HM}_{\max}^{\text{direct}}]$ — this normalises each
          panel to its own maximum and directly compares the spatial
          distribution patterns regardless of magnitude.
\end{itemize}

\paragraph{Interpretation.}
A discrepancy in absolute stress magnitude between the coupled and direct
models is expected: the frame model predicts a different (generally smaller)
tip displacement than the true solid, because Euler--Bernoulli theory
overestimates stiffness for the non-slender geometry used here
($R/l_s \approx 0.93$, see Section~\ref{sec:locking}).  The independent-scale
figures therefore provide the more informative comparison: if the normalised
hot-spot patterns are spatially similar, the frame-based approach captures the
correct stress topology even if it under- or over-estimates the magnitude.

\begin{figure}[htbp]
\centering
\begin{minipage}{0.48\textwidth}
\centering
\exampleimage{coupledExample_sampleMinN_smooth_deformed.png}{Run \texttt{coupledExamples.m} to generate\\\texttt{coupledExample\_sampleMinN\_smooth\_deformed.png}}
\end{minipage}\hfill
\begin{minipage}{0.48\textwidth}
\centering
\exampleimage{coupledExample_sampleMinN_smooth_stress.png}{Run \texttt{coupledExamples.m} to generate\\\texttt{coupledExample\_sampleMinN\_smooth\_stress.png}}
\end{minipage}
\caption{Configuration \texttt{sampleMinN\_smooth} = $[0,180,180,180,180,180,180]$. Left: deformed mesh comparison. Right: Huber--Mises stress, \textbf{shared} color scale.}
\end{figure}

\begin{figure}[htbp]
\centering
\exampleimage{coupledExample_sampleMinN_smooth_stress_normalized.png}{Run \texttt{coupledExamples.m} to generate\\\texttt{coupledExample\_sampleMinN\_smooth\_stress\_normalized.png}}
\caption{Configuration \texttt{sampleMinN\_smooth}: Huber--Mises stress with \textbf{independent} per-panel color scales.  Spatial hot-spot patterns can be compared directly regardless of magnitude.}
\end{figure}

\begin{figure}[htbp]
\centering
\begin{minipage}{0.48\textwidth}
\centering
\exampleimage{coupledExample_sampleMaxMz_smooth_deformed.png}{Run \texttt{coupledExamples.m} to generate\\\texttt{coupledExample\_sampleMaxMz\_smooth\_deformed.png}}
\end{minipage}\hfill
\begin{minipage}{0.48\textwidth}
\centering
\exampleimage{coupledExample_sampleMaxMz_smooth_stress.png}{Run \texttt{coupledExamples.m} to generate\\\texttt{coupledExample\_sampleMaxMz\_smooth\_stress.png}}
\end{minipage}
\caption{Configuration \texttt{sampleMaxMz\_smooth} = $[0,0,0,180,180,180,180]$. Left: deformed mesh comparison. Right: Huber--Mises stress, \textbf{shared} color scale.}
\end{figure}

\begin{figure}[htbp]
\centering
\exampleimage{coupledExample_sampleMaxMz_smooth_stress_normalized.png}{Run \texttt{coupledExamples.m} to generate\\\texttt{coupledExample\_sampleMaxMz\_smooth\_stress\_normalized.png}}
\caption{Configuration \texttt{sampleMaxMz\_smooth}: Huber--Mises stress with \textbf{independent} per-panel color scales.}
\end{figure}

\begin{figure}[htbp]
\centering
\begin{minipage}{0.48\textwidth}
\centering
\exampleimage{coupledExample_sampleMaxTy_smooth_deformed.png}{Run \texttt{coupledExamples.m} to generate\\\texttt{coupledExample\_sampleMaxTy\_smooth\_deformed.png}}
\end{minipage}\hfill
\begin{minipage}{0.48\textwidth}
\centering
\exampleimage{coupledExample_sampleMaxTy_smooth_stress.png}{Run \texttt{coupledExamples.m} to generate\\\texttt{coupledExample\_sampleMaxTy\_smooth\_stress.png}}
\end{minipage}
\caption{Configuration \texttt{sampleMaxTy\_smooth} = $[0,0,180,0,180,180,180]$. Left: deformed mesh comparison. Right: Huber--Mises stress, \textbf{shared} color scale.}
\end{figure}

\begin{figure}[htbp]
\centering
\exampleimage{coupledExample_sampleMaxTy_smooth_stress_normalized.png}{Run \texttt{coupledExamples.m} to generate\\\texttt{coupledExample\_sampleMaxTy\_smooth\_stress\_normalized.png}}
\caption{Configuration \texttt{sampleMaxTy\_smooth}: Huber--Mises stress with \textbf{independent} per-panel color scales.}
\end{figure}

\begin{figure}[htbp]
\centering
\begin{minipage}{0.48\textwidth}
\centering
\exampleimage{coupledExample_sampleMaxMs_smooth_deformed.png}{Run \texttt{coupledExamples.m} to generate\\\texttt{coupledExample\_sampleMaxMs\_smooth\_deformed.png}}
\end{minipage}\hfill
\begin{minipage}{0.48\textwidth}
\centering
\exampleimage{coupledExample_sampleMaxMs_smooth_stress.png}{Run \texttt{coupledExamples.m} to generate\\\texttt{coupledExample\_sampleMaxMs\_smooth\_stress.png}}
\end{minipage}
\caption{Configuration \texttt{sampleMaxMs\_smooth} = $[0,45,45,45,270,180,180]$. Left: deformed mesh comparison. Right: Huber--Mises stress, \textbf{shared} color scale.}
\end{figure}

\begin{figure}[htbp]
\centering
\exampleimage{coupledExample_sampleMaxMs_smooth_stress_normalized.png}{Run \texttt{coupledExamples.m} to generate\\\texttt{coupledExample\_sampleMaxMs\_smooth\_stress\_normalized.png}}
\caption{Configuration \texttt{sampleMaxMs\_smooth}: Huber--Mises stress with \textbf{independent} per-panel color scales.}
\end{figure}

\end{document}
